% Compile with XeLaTeX. This file is standalone; references are mapped in section 4 of the package.
\documentclass[UTF8,a4paper,11pt,fontset=fandol]{ctexart}
\usepackage[left=23mm,right=23mm,top=25mm,bottom=24mm,headheight=16pt]{geometry}
\usepackage{fontspec}
\setmainfont{lmroman10-regular.otf}[BoldFont=lmroman10-bold.otf,ItalicFont=lmroman10-italic.otf,BoldItalicFont=lmroman10-bolditalic.otf]
\setsansfont{lmsans10-regular.otf}[BoldFont=lmsans10-bold.otf,ItalicFont=lmsans10-oblique.otf,BoldItalicFont=lmsans10-boldoblique.otf]
\setmonofont{lmmono10-regular.otf}[Scale=0.84,ItalicFont=lmmono10-italic.otf,BoldFont=lmmonolt10-bold.otf]
\usepackage{amsmath,amssymb,mathtools,bm,esint,mathrsfs}
\usepackage{longtable,booktabs,array,calc,etoolbox}
\usepackage{graphicx}
\usepackage{fancyvrb}
\usepackage{xurl}
\usepackage{fancyhdr,lastpage}
\usepackage[unicode,hidelinks]{hyperref}
\hypersetup{pdftitle={数学深化：关键定理与攻关路线},pdfauthor={Qiushi Engine research handoff},pdfsubject={Quantitative Einstein compatibility and mathematical physics}}
\setlength{\parindent}{0pt}
\setlength{\parskip}{0.45em}
\setlength{\emergencystretch}{3em}
\setlength{\tabcolsep}{5pt}
\renewcommand{\arraystretch}{1.22}
\AtBeginEnvironment{longtable}{\small}
\setcounter{secnumdepth}{0}
\setcounter{tocdepth}{2}
\ctexset{
 section={format=\Large\bfseries\raggedright,beforeskip=1.3em,afterskip=0.6em},
 subsection={format=\large\bfseries\raggedright,beforeskip=1em,afterskip=0.4em},
 subsubsection={format=\normalsize\bfseries\raggedright,beforeskip=0.8em,afterskip=0.3em}}
\pagestyle{fancy}
\fancyhf{}
\fancyhead[L]{\small\sffamily Qiushi Engine · 第一科学主线}
\fancyhead[R]{\small 2026-09-28}
\fancyfoot[C]{\small \thepage\ / \pageref{LastPage}}
\renewcommand{\headrulewidth}{0.3pt}
\renewcommand{\footrulewidth}{0pt}
\providecommand{\tightlist}{\setlength{\itemsep}{0pt}\setlength{\parskip}{0pt}}
\providecommand{\pandocbounded}[1]{#1}
\newcommand{\handoffpart}[1]{\clearpage\section{#1}}
\begin{document}
\begin{titlepage}
\thispagestyle{empty}
\vspace*{2cm}
{\large\sffamily QIUSHI ENGINE\par}
\vspace{1.3cm}
{\Huge\bfseries 数学深化：关键定理与攻关路线\par}
\vspace{0.8cm}
{\Large 从线性强制估计到缺陷传递、最佳首项与全局实现\par}
\vspace{1.2cm}
{\large 跨会话研究综合与证明方案\par}
\vfill
{\normalsize 原稿状态基线：2026年9月27日，快照 \texttt{844e62f}\par}
\vspace{0.35cm}
{\normalsize 整理日期：2026年9月28日\par}
\vspace{1cm}
{\small 保留已取得的强结论，沿真实方程继续推进全局数学与物理问题。\par}
\end{titlepage}
\tableofcontents
\clearpage
\section{1.
从已经取得的最强状态出发}\label{ux4eceux5df2ux7ecfux53d6ux5f97ux7684ux6700ux5f3aux72b6ux6001ux51faux53d1}

本文件承接 \texttt{844e62f}
原稿包及9月27日统一任务书，不把早稿尚未完成的陈述继续当作最新状态。{[}S0--S2{]}

\begin{longtable}[]{@{}
  >{\raggedright\arraybackslash}p{(\columnwidth - 4\tabcolsep) * \real{0.3333}}
  >{\raggedright\arraybackslash}p{(\columnwidth - 4\tabcolsep) * \real{0.3333}}
  >{\raggedright\arraybackslash}p{(\columnwidth - 4\tabcolsep) * \real{0.3333}}@{}}
\toprule\noalign{}
\begin{minipage}[b]{\linewidth}\raggedright
项目
\end{minipage} & \begin{minipage}[b]{\linewidth}\raggedright
当前状态
\end{minipage} & \begin{minipage}[b]{\linewidth}\raggedright
下一项有实质价值的推进
\end{minipage} \\
\midrule\noalign{}
\endhead
\bottomrule\noalign{}
\endlastfoot
I40：小 \(h\) 线性阶 & 指定真实平衡族已闭合 & 最佳首项与近极小几何 \\
I41：无标量方差强制估计 & 已闭合，全度量适用 &
自然输入扩展、真正新增的全局后果 \\
BF 新线性窗口 & 已有实际类供给与拓扑构造 &
完整经典预算比较、更多自然实例 \\
I44：只添柄子族 & 已有精确 \(\inf\mathcal E\) &
与混合平衡族分开，作为锐利标定 \\
I45：有界 Ricci 与 Euler 密度零端 & 已闭合；指定逐点输入族已加强 &
删除点态 Ricci 控制，再研究删除密度比值 \\
I47：完整光滑核心 & 指定取向、锥领口下代价锐利 &
从原弱输入提取带标记核心并接回全局 \\
I48：椭圆／驻点接口 & 各自有附加条件的已证结果 &
证明原目标序列满足接口，或发展不同方法 \\
I49：共形 Ricci 流 & 实际近极小初值的时间上界已给出 &
如需推广，研究真实演化方程及奇异结构 \\
X06：Weyl 支付系数 & 已有真实反例，普遍删去该项不可行 &
利用平衡族信息控制保留该项后的有符号余量 \\
I25/I26 与覆盖实现 & 保留各自输入范围内的精确实现 &
自然边界去直径、统一修正和全局标记 \\
固定 \(E_D\) & 已有逐点近负 Einstein 构造；全部 Einstein 判定未闭合 &
覆盖所有候选度量的存在或排除机制 \\
\texttt{e58b867} 高指数模型 & 本套材料中为摘要接续 &
读取完整推导，完成全局适配与代价判定 \\
\end{longtable}

目前最能提高理论深度的共同问题是：\textbf{实际方程如何控制那些局部可以隐藏、但全局不能免费消失的缺陷。}首项问题与零端问题可共享分析，但输出精度不同。

\section{2.
固定记号与实际供给合同}\label{ux56faux5b9aux8bb0ux53f7ux4e0eux5b9eux9645ux4f9bux7ed9ux5408ux540c}

采用闭定向四维黎曼流形、原稿曲率范数和取向。令

\[
\begin{gathered}
V=\operatorname{Vol}(M,g),\quad
P=\int|\operatorname{Ric}^{\circ}|^2,\quad R=\int|\operatorname{Ric}|^2,\\
T_\pm=2\chi\pm3\tau,\quad p=(a_g^+)^2,\quad
\Sigma=\int s^2-32\pi^2p,\\
D_2=\frac32\int s^2+18\int|W^+|^2-72\pi^2p,
\quad Y=(T_+-a^2)_+.
\end{gathered}
\]

积分均使用 \(dV_g\)。\(a^2\) 是拓扑交叉平方；\(p\) 是随度量正 Hodge
子空间变化的投影平方。\textbf{不能把 \(p\) 视为固定拓扑常数。}预算
\(A>0\) 是逐度量实际供给所保证的统一下界 \(p\ge A\)。

``实际供给''的意思不是代数上存在某个陈类，而是所选真实结构在证明所需共形／权重范围内有真实规范解。同一初始度量下不能为每条估计挑不同结构，随后把它们当成同一个场相加。BF
与 Pin
型供给属于不同拓扑入口；如果最终产生同类方程，推广主要体现为适用对象扩展，而非自动产生第二套分析机制。

在含多个记号 \(A\) 的正式写作中，建议将联络写为 \(\mathcal A\)，保留
\(A\) 给拓扑预算。本文件以下采用这一约定。

\section{3.
核心证明为什么闭合：可供迁移的技术骨架}\label{ux6838ux5fc3ux8bc1ux660eux4e3aux4ec0ux4e48ux95edux5408ux53efux4f9bux8fc1ux79fbux7684ux6280ux672fux9aa8ux67b6}

I40 的关键量为

\[
K=\int|\nabla_{\mathcal A}\Phi|^2,\quad
B=\int|F_{\mathcal A}^-|^2,\quad
\Theta=rV+D_2+\varepsilon V,
\quad r=\fint|\operatorname{Ric}+g|^2.
\]

校准幅度、旋量范数、Weyl 模型与权重全部来自同一实际解。设
\(H=\nabla_{\mathcal A}^2\Phi\)、\(L=\operatorname{tr}H\)、\(P_i=\nabla_i\Phi\)。核心
Hessian 恒等式保留曲率的作用结构：

\[
\begin{aligned}
\|H\|_2^2={}&\|L\|_2^2
-\int\operatorname{Ric}_{ij}\operatorname{Re}\langle P_i,P_j\rangle\\
&+\sum_{i,j}\int\operatorname{Re}\langle F^{S^+}_{ij}P_j,P_i\rangle
+\frac12\sum_{i,j}\|F^{S^+}_{ij}\Phi\|_2^2.
\end{aligned}
\tag{M1}
\]

这里的双重求和规范与外形式曲率范数可能相差固定因子，应沿原稿使用，不能混换。式(M1)不对曲率或权重求导。{[}S0，卷01
Source 1，§4.1{]}

点态代数背景满足

\[
|F_{\mathrm{bg}}|\le C,\qquad
|F_{\mathrm{bg}}\Phi|^2\le C(\rho-\rho_0)^2.
\]

因此大的背景只支付 \(CK\)，而不再支付独立的
\(CV\)。结合有界截面分部积分和实际旋量方程，得到

\[
\|H\|_2^2+\|\nabla_{\mathcal A}\Phi\|_4^4
\le C(K+\Theta+B).
\tag{M2}
\]

另一方面，同解行列式能量与曲率恒等式给出

\[
B\le\sqrt6K+C\Theta
\tag{M3}
\]

（此处使用 \(a^2\ge T_+\)）。弱 Bianchi 测试再给

\[
K\le C\Theta+C\sqrt{\Theta K},
\quad\text{从而}\quad K\le C\Theta.
\tag{M4}
\]

低幅度区域由方程的截断测试控制，不通过除以 \(|\Phi|\) 得到。I41
对归一化结果做精确共形转移，最终得到

\[
\Sigma\le C\left(P+D_2+\sqrt{(P+D_2)Y}\right).
\tag{M5}
\]

这一技术骨架可以提炼为一般证明模板，但只有在新方程中逐项证明实际供给、背景消去、拓扑吸收和自测试，才构成真正推广。把这些条件重新命名后列为假设，不等于已经建立了新的物理理论。

\section{4.
两个立即有用的代数接口}\label{ux4e24ux4e2aux7acbux5373ux6709ux7528ux7684ux4ee3ux6570ux63a5ux53e3}

\subsection{4.1
拓扑偏离强迫总缺陷预算}\label{ux62d3ux6251ux504fux79bbux5f3aux8febux603bux7f3aux9677ux9884ux7b97}

原稿的精确恒等式为

\[
36\pi^2(p-T_+)=\frac92P+\frac98\Sigma-D_2.
\tag{M6}
\]

在 \(p\ge A\)、\(T_+=A(1-h)\)、\(h>0\)、\(a^2\ge T_+\) 下，令

\[
K_C=\frac92+\frac98C,
\]

其中 \(C\) 来自 \(\Sigma\le C(P+D_2)\)。直接有

\[
36\pi^2Ah\le\frac92P+\frac98\Sigma-D_2
\le K_C(P+D_2).
\]

所以

\[
\boxed{P+D_2\ge\frac{36\pi^2}{K_C}Ah.}
\tag{M7}
\]

这是本套文件补充写出的直接推论。它不需要把 \(U\)
预先放在平衡窗口，适合先连接正系数的二次曲率作用量。要只由 \(P\)
支付代价，则仍需式(M8)所用的双手性信息。

\subsection{4.2 平衡区间强迫无迹 Ricci
本身付费}\label{ux5e73ux8861ux533aux95f4ux5f3aux8febux65e0ux8ff9-ricci-ux672cux8eabux4ed8ux8d39}

新线性窗口是

\[
\frac{4P}{R}+(U-1-h)_+\ge ch.
\tag{M8}
\]

若所选真实族满足 \((U-1-h)_+\le(c/2)h\)，则

\[
\frac{4P}{R}\ge\frac c2h.
\]

再使用该实际输入下的经典预算

\[
R\ge8\pi^2(2A-T_+)=8\pi^2A(1+h),
\]

得到

\[
\boxed{P\ge C_0Ah}
\tag{M9}
\]

（可取相应的固定
\(C_0>0\)）。式(M9)是物理中最小无迹应力下界的入口。它不是无条件的``任意拓扑荷都产生引力能量''。{[}S0，卷01；R1、R4{]}

\subsection{4.3
动态推论的准确位置}\label{ux52a8ux6001ux63a8ux8bbaux7684ux51c6ux786eux4f4dux7f6e}

I49 已将三项静态预算送入 \(s=-4\) 归一化的共形 Ricci 流。对真实平衡族
\(h_j\to0\)，选取满足

\[
\mathcal E(g_j(0))\le\frac{4h_j}{1+h_j}+h_j^2
\]

的实际初值，若 \(0<c_1<4\) 且 \(c_1h_j\) 是已知全度量下界，原稿得到

\[
\limsup_{j\to\infty}T_{\max,j}\le1+\log(4/c_1).
\tag{M9a}
\]

它不要求极小值达到，显示静态强制估计已经能够影响一种真实几何演化。其对象是原稿的共形
Ricci 流及其压力规范，不是普通归一化 Ricci
流，更不是物理宇宙的演化时间。{[}S0，导读§4.4；卷03 Source 2，§1{]}

\section{5.
首项问题：决定代价，而不是继续泛化``更锐利''}\label{ux9996ux9879ux95eeux9898ux51b3ux5b9aux4ee3ux4ef7ux800cux4e0dux662fux7ee7ux7eedux6cdbux5316ux66f4ux9510ux5229}

固定原稿中同一个真实平衡族，设

\[
\alpha_j=\inf_{g\in\operatorname{Met}(M_j)}\mathcal E(g),\qquad
c_1h_j\le\alpha_j\le\frac{4h_j}{1+h_j}.
\]

\textbf{目标 A：决定 \(\alpha_j/h_j\)
的极限或上下极限，并解释实现该值的几何。}允许答案为4，也允许同族新构造给出更低值；不能从经典拼接的漂亮形式预设其全局最优。

I41 的共形约化允许在实际常标量 \(s=-4\)
的近极小序列中研究。首项为4可接到

\[
\Sigma\le\frac89D_2+32\pi^2(p-A)+o(Ah),
\tag{M10}
\]

或等价的有符号余量目标

\[
P-8\pi^2Ah=2(2F-\mathscr S),
\qquad
\boxed{(\mathscr S-2F)_+=o(Ah)}.
\tag{M11}
\]

此处 \(F=\int|W^+-W_0|^2\)；\(\mathscr S\) 必须是原稿同一实际 SW
解给出的完整有符号余量，含幅度、相位、梯度与投影超额，不能另选一个独立标量。{[}S1，§四.1；S2，§6{]}

现在有三条明确证据指导取舍。其一，单纯优化 I41
各向同性常数不足以逼出4，真实乘积模型已经限制该方法。其二，X06 给
\(\mathscr S/F\to2\)、\(P/F\to0\)，否定普遍删去 \(2F\) 的企图；但该例在
\(h=0,U=3\)，不否定平衡族中的式(M11)。其三，纯共形 Hessian
已为正，若寻找光滑临界点的负方向，应研究约化 TT
部分，而不是继续尝试纯共形扰动。{[}S0，卷03 Source 2，§§3--4{]}

最值得探索的是：另一手性的紧预算怎样限制 \(\mathscr S\)
可以利用的区域与尺度；高指数过渡层是否能够在 \(Ah\)
尺度隐藏负余量；新的真实拼接能否通过耦合而低于经典4。每个局部估计都应返回式(M11)，检查它是否真的把剩余量压到
\(o(Ah)\)。仅有 \(o(V)\) 不够。

只添柄复双曲子族的精确结果

\[
\inf_g\mathcal E(g)=\frac{4h}{1+h},\qquad0\le h\le1/3,
\]

是重要标定，但不能替代混合吹起／添柄平衡族的首项判定。{[}S0，卷02{]}

\section{6.
一般弱积分零端：当前最准确的方程目标}\label{ux4e00ux822cux5f31ux79efux5206ux96f6ux7aefux5f53ux524dux6700ux51c6ux786eux7684ux65b9ux7a0bux76eeux6807}

对 \(\beta\ge2\)，设

\[
D_\beta=\frac{\beta+1}{\beta}
\int(s^2+6\beta|W^+|^2)-72\pi^2p,
\qquad\kappa_\beta=r+\frac{D_\beta}{V}.
\]

\textbf{目标 B：在原实际加权方程及类平方条件下，建立}

\[
\boxed{\kappa_\beta\ge c\frac{(\beta-2)^2}{\beta^2},
\qquad2\le\beta<\infty.}
\tag{M12}
\]

常数不额外依赖点态 Ricci 界、直径、非塌缩、统一 Sobolev
常数或驻点方程。此式是目标，不是当前一般定理。

更接近现有证明且立即有真实应用的台阶是：保留

\[
p\ge A>0,\quad a^2\ge T_+,\quad\chi(M)/A\le B_0,
\]

完全删除 I45 的点态 Ricci 控制，取得常数 \(c(B_0)>0\)
的纯积分版本。原稿指定的一般型手术族已有相应拓扑比例控制，因此这一步能直接扩大实际应用，而不只是形式删除假设。{[}S2，§3{]}

已有紧参数区间估计把困难压到假想序列

\[
\beta_j\to\infty,\qquad\kappa_{\beta_j}\to0.
\tag{M13}
\]

这个零端与 KE 端不是同一种近几何。原方程的归一化强迫

\[
\fint|\nabla\omega|^2\to\frac83,
\tag{M14}
\]

所以不能套用``形式趋近平行''的证明。实际扭率低体积过渡层则必须保留正的临界混合密度：

\[
\frac1V\int_{T_\beta}
|\langle t,\nabla t\rangle|^2|\nabla t|^2\ge c>0.
\tag{M15}
\]

可能的突破有两种。第一，从完整实际方程证明式(M15)左侧不可能保持正量，得到全局矛盾。第二，完成满足全部原始输入的真实逃逸序列，确定所希望统一间隙的准确失效范围。两者都比不断添加没有改变应用对象的条件版本更有价值。

I45、I48
仍然是工具库中的强结论。但任意近极小序列不自动满足二次曲率驻点方程，平均
Ricci 小误差也不自动给出局部 \(L^p\)
小性和覆盖正则性，不能将它们未经连接直接用于式(M13)。

\section{7.
高指数零电荷模型：寻找遗漏的全局约束}\label{ux9ad8ux6307ux6570ux96f6ux7535ux8377ux6a21ux578bux5bfbux627eux9057ux6f0fux7684ux5168ux5c40ux7ea6ux675f}

\texttt{e58b867} 的摘要报告真实局部规范解具有扭率指数
\(d^2\)，而相对行列式电荷为零。其意义不是已经得到闭流形反例，而是：\textbf{一种电荷守恒不足以控制所有实际场复杂度。}{[}S2，§4{]}

首先读取完整方程、权重和边界平凡化。随后将指数 \(d\)、插入尺度
\(r\)、外部体积和拓扑预算同时作为变量，完成以下适配表。

\begin{longtable}[]{@{}
  >{\raggedright\arraybackslash}p{(\columnwidth - 2\tabcolsep) * \real{0.5000}}
  >{\raggedright\arraybackslash}p{(\columnwidth - 2\tabcolsep) * \real{0.5000}}@{}}
\toprule\noalign{}
\begin{minipage}[b]{\linewidth}\raggedright
对象
\end{minipage} & \begin{minipage}[b]{\linewidth}\raggedright
需要实际计算的内容
\end{minipage} \\
\midrule\noalign{}
\endhead
\bottomrule\noalign{}
\endlastfoot
局部规范系统 & 满足原系统哪些方程，是否改变权重或扰动 \\
旋量与二形式 & 幅度、方向、相位在零端标定下怎样变化 \\
临界预算 &
\(\int|\nabla\Phi|^4\)、\(\int|\nabla^2\Phi|^2\)、扭率混合项随 \(d,r\)
的依赖 \\
几何预算 & Ricci、两手性 Weyl、标量松弛及 Euler 比例能否同时受控 \\
边界与传递 & 行列式、秩二特征量、holonomy 与有符号通量怎样匹配 \\
全局闭合 & 外部区域与过渡带需支付的最小代价 \\
供给量词 & 得到一个全局解，是否仍满足原命题所需整个共形／权重范围 \\
\end{longtable}

四维临界缩放提示：不能仅靠把插入区域缩小，假定所有预算随体积消失。需要直接算出真正不随尺度衰减的项。

优先从现有秩二联络的 Chern--Weil／transgression
恒等式辨认缺失信息。若需要新的标记，应说明它如何由实际场定义、在边界上怎样变化、凭什么得到强制下界；不要先定义一个没有控制的``总电荷''，再假定它阻止模型。

成功的输出应是一个完整连接：高指数模型在全局必须付出 \(V\) 或 \(Ah\)
尺度的代价，且覆盖目标序列；或一个真正满足原假设的闭流形构造。继续展示更多局部指数值，不等同于完成全局任务。

\section{8.
完整核心、相对特征量与跨尺度守恒}\label{ux5b8cux6574ux6838ux5fc3ux76f8ux5bf9ux7279ux5f81ux91cfux4e0eux8de8ux5c3aux5ea6ux5b88ux6052}

对 I47 的商阶 \(n\)，令

\[
\mathscr Q=|\operatorname{Ric}^{\circ}|^2+4|W^+|^2+s^2/12,
\qquad B_n=3-\frac n2-\frac1{2n}.
\]

完整锥领口填充满足

\[
2\chi(\Omega)+3\tau(\Omega)-B_n
=\frac1{4\pi^2}\int_\Omega
\left(2|W^+|^2+\frac{s^2}{24}-\frac{|\operatorname{Ric}^{\circ}|^2}{2}\right).
\tag{M16}
\]

由于

\[
\operatorname{dist}(B_n,\mathbb Z)=\frac{n-1}{2n},
\]

式(M16)给 \(\int\mathscr Q\ge4\pi^2(n-1)/n\)。经典 ALE
截断与领口展平提供匹配上界。{[}S0，卷03 Source 1，§1{]}

这条结论依赖取向和完整领口，不是``任意 lens space
边界都支付同一个能量''。例如反转取向会改变允许的 Ricci-flat
填充。若领口仅近锥，原稿已给带明确误差的版本。

一个中间区域若有一个外边界和多个内边界，相对项是

\[
-B_{n_0}+\sum_iB_{n_i},
\]

内部边界在拼接后抵消。因此必须对整个标记核心计费，不能对泡树每个节点重复计费。守恒公式与非负强制代价是两步不同的证明。

当前尚需解决的是三个连接：原弱输入是否产生适用区域及受控重叠；非零实际场是否在同一泡尺度存活；完整手术区的标记如何从外部传递至内部并汇总。全局误差趋零不推出其除以
\(r_j^4\) 后仍小，也不能假定有限 ALE 泡树已经存在。{[}S1，§四.3{]}

\section{9.
有效缺陷理论：应证明什么，而非预设什么}\label{ux6709ux6548ux7f3aux9677ux7406ux8bbaux5e94ux8bc1ux660eux4ec0ux4e48ux800cux975eux9884ux8bbeux4ec0ux4e48}

``近 Einstein
几何的拓扑缺陷有效理论''是有价值的研究纲领：它试图解释宏观线性代价怎样由局部结构及其耦合产生。当前尚不是一项完成的
\(\Gamma\) 极限定理。

固定实际族 \(M_j\)、\(h_j\to0\)，考虑

\[
\mathcal F_j(g)=\frac{\mathcal E(g)}{h_j}.
\]

要谈变分极限，必须先规定变动流形上的状态对象与收敛：度量、实际场、边界平凡化、相对特征量、规范等价，以及保留下来的尺度。随后证明足够的紧性、下极限不等式和恢复序列。{[}R12{]}

有意义的输出是

\[
\mathcal F_{\mathrm{eff}}(\text{标记缺陷状态}),
\]

其具体形式由证明决定。可以包含核心、颈部、单值变换和弥散部分，也可能存在耦合和抵消。\textbf{不能预设它就是独立缺陷能量之和。}I47
的单核心代价既没有确定平衡族的全部核心类型，也没有给出全部缺陷的加法定律。

较短的路线同样值得优先：若直接证明式(M11)并给恢复序列，已经可以决定首项，不必等待完整有效理论。反之，一个确实覆盖原输入的标记分解，若同时解释首项和零端，会有更深的统一价值。两条路线应以实际证明进展选择，而不是以抽象程度高低选择。

\section{\texorpdfstring{10. 固定 \(E_D\)：直接触及 Yau 26
的全局检验}{10. 固定 E\_D：直接触及 Yau 26 的全局检验}}\label{ux56faux5b9a-e_dux76f4ux63a5ux89e6ux53ca-yau-26-ux7684ux5168ux5c40ux68c0ux9a8c}

取纤维和基底均为亏格2曲面的丛；纤维上非分离曲线 \(\alpha\) 的单值变换为

\[
\rho(a_1)=D_\alpha,\qquad
\rho(b_1)=\rho(a_2)=\rho(b_2)=1.
\]

现有深化报告接续的原稿构造为

\[
\chi(E_D)=4,\quad\tau(E_D)=0,\quad b_1(E_D)=7,
\]

以及同一光滑流形上的

\[
\delta_-(g_j)\to0,\qquad V(g_j)=16\pi^2.
\tag{M17}
\]

式(M17)是逐点逼近构造，不能被变动流形族的积分小量替代。本套文件按{[}S3，§9{]}接续其状态；接收会话在使用反例结论前，应直接调用卷08的原始构造及核查。

\textbf{目标 C：决定 \(E_D\) 是否存在任何满足 \(\operatorname{Ric}=-g\)
的度量。}

\(b_1=7\) 通过 Bochner 型约束排除正 Einstein 和 Ricci-flat 情形；负
Einstein 的全部候选仍留有体积窗口

\[
16\pi^2<V(g_E)<48\pi^2.
\tag{M18}
\]

当前规范预算为 \(A=T_+=T_-=8\)，故
\(h=0,U=3\)。它不位于新线性壁最强的平衡区间。优化 I41
的普遍常数，不能自动解决这个固定对象。

若排除全部负 Einstein 候选，式(M17)即构成项目所研究纯 Ricci
夹逼命题的负号分支反例。原始 Yau 26
的符号量词仍须单独表述：否定负号分支，不等于否定存在某一正号常数窗口的整个命题。固定维数下变动流形的严格反例序列也可能足以否定维数统一的负号版本，并非必须固定
\(E_D\)。

现有普通单极子凸包预算、标准熵判据和某些现成填充定理均不能直接关闭式(M18)。值得尝试的新机制，是从完整
Einstein 方程和曲面丛群结构导出对无限 Dehn twist
的限制，或构造不要求输出接近 \(g_j\)
的新精确度量。调和映射、最小曲面、模空间慢场及全局周期可以参与，但必须证明相应结构由任意候选
Einstein 度量产生，不能预设它保纤维。

只证明原逼近序列附近无解、某种算法失败、能量下确界不达到，都没有覆盖式(M18)中的全部度量。

\section{11.
自然边界的精确实现与慢模态}\label{ux81eaux7136ux8fb9ux754cux7684ux7cbeux786eux5b9eux73b0ux4e0eux6162ux6a21ux6001}

另一个重大目标位于

\[
\chi=3\tau>0,\qquad T_+=3T_-.
\]

在实际供给 \(p\ge T_+\) 和拓扑上界 \(T_+\le A_{\max}\)
下，争取从足够小的 \(\delta_-\)
推出同一光滑流形承载复双曲度量，而不预设直径。已有恒等式已把体积和多种曲率缺陷接到拓扑预算上；剩下的是有限体积长颈、塌缩及端部填充是否允许逃逸。{[}S3，§7{]}

这与 \(\tau=0\)
的双圆盘剪切族不同，不能把后者的反例直接移到此边界。较近的可判定结果，是在原有界直径范围证明输入度量到复双曲代表的定量修正速率；需要一致的规范切片和线性逆算子控制，而不是仅有序列紧性。

慢模态分支则要保留完整 Teichmüller 值场、基底核和连接变量。I33
的全张量小邻域与 I36
的固定幅度特殊模型不能直接合成``固定幅度、覆盖次数无关、全部 Einstein
输出''的定理。若要与紧化物理连接，应先证明法向方向一致可逆，随后求解所有低能投影方程。关键是完整非线性反馈，不是单独记录逆算子增长率。{[}S3，§10{]}

\section{12.
多实际场与矩阵证书：独立的可检验入口}\label{ux591aux5b9eux9645ux573aux4e0eux77e9ux9635ux8bc1ux4e66ux72ecux7acbux7684ux53efux68c0ux9a8cux5165ux53e3}

原材料中还保留一个无需先固定全局相对符号的联合判据。令

\[
Q_{ij}=\frac{2\pi^2}{V}(a_i\cdot a_j),\quad
H(Z)=\max_{\sigma\in\{\pm1\}^m}\sigma^{\mathsf T}Z\sigma,
\quad Z\succeq0.
\]

在原 I12 输入及小缺陷范围中，辅助推导给

\[
\operatorname{tr}(ZQ)
\le\left(\sqrt{H(Z)}+
\sqrt{Cm\kappa\operatorname{tr}Z}\right)^2.
\tag{M19}
\]

零误差极限因此受

\[
Q\in\operatorname{conv}\{\sigma\sigma^{\mathsf T}\}
-\operatorname{PSD}_m
\]

约束。有限维半正定分离可给精确证书，但实际单极子类配置必须先被构造出来。标准正交
BF
符号块可能天然处于允许区，不能靠反复优化同一矩阵期待新实例。{[}S3，§6及附录A{]}

它适合作为一条有清晰停止条件的支线：要么得到真实分离族，要么证明所选供给永不分离，从而知道必须引入新的微分配对或全局周期信息。最新
BF 线性窗口来自 I41，并不是式(M19)已经找到了分离族。

\section{13.
扩展为一般相容性原理的最低要求}\label{ux6269ux5c55ux4e3aux4e00ux822cux76f8ux5bb9ux6027ux539fux7406ux7684ux6700ux4f4eux8981ux6c42}

可将现有技术组织为四个接口：实际解供给、背景作用消去、剩余能量的拓扑控制、微分测试闭合。一般化的价值取决于新实例，而不是把所有困难装入抽象假设。

一个值得建立的结果应当同时做到：原有 SW/BF/Pin
应用以准确输入进入；至少一个自然新方程或新几何类别得到原来没有的全局结论；最坏退化由新机制控制；必要假设的失效有实际模型解释。

在这种意义上，研究可以进入新的规范体系、不同校准几何、具有边界的相对问题或耦合物理系统。不能预设所有相容几何最终都
Kähler，也不能预设高维存在相同的四维手性分解。可迁移的是作用消去和相容性估计，具体表示与特征类需要重建。

\section{14.
两个接收会话如何协同推进}\label{ux4e24ux4e2aux63a5ux6536ux4f1aux8bddux5982ux4f55ux534fux540cux63a8ux8fdb}

两个会话保持同一份当前状态、同一套符号和一个选定的全局目标。建议一方主攻实际方程与有符号估计，另一方主攻真实构造、最坏局部模型和全局匹配；两方共享同一个假设合同，而不是分别发展无法比较的版本。

下一轮优先读取高指数模型的完整提交，并把它代入零端或首项所需的准确预算。若某条直接方程路线比完整分解更短，应立即采用。若某个局部模型迫使改变研究方向，应记录它否定的确切命题，继续研究仍可能成立的全局结论。

每一阶段的实质交付只有三类：一个带完整证明的新命题；一个满足相同假设的真实反例；一个经计算与文献检验、比此前更短更准确的剩余连接。符号脚本、区间计算、精确矩阵证书、有限元或高精度数值都可服务这些输出，但不替代无限维定理的证明。

原有主论文同步整理，不与继续攻关互相等待。研究价值的最终尺度，是新机制让哪一种过去无法判断的几何变得可判断。
\end{document}
